\documentclass[11pt]{article}

\usepackage{amsmath,amssymb}
\usepackage[margin=1in]{geometry}
\usepackage{booktabs}
\usepackage{array}
\usepackage{hyperref}

\title{Ordering, Causal Closure, and Proto-Geometric Distance-Surrogate Feasibility on a Frozen Branch-Local Cochain-Laplacian Hierarchy}
\author{Tomislav Rupi\'c \\ Independent Researcher}
\date{March 2026}

\begin{document}

\maketitle

\begin{abstract}
This synthesis paper freezes the Phase XVI--XVIII regime of the HAOS-IIP program as one bounded kinematics stack. The question is not whether spacetime or metric geometry emerges. The question is narrower: on the already frozen branch-local cochain-Laplacian hierarchy $\Delta_h$, does there exist a refinement-stable chain
\[
\text{propagation ordering}
\;\to\;
\text{causal closure}
\;\to\;
\text{distance-surrogate compatibility}
\]
that survives matched deterministic controls? Phase XVI established a stable event-ordering structure using collective observables alone, with branch primary event distance $0.142857142857$, primary ordering score $0.965986394558$, front-order distance $0.192592592593$, and robustness distance $0.01421957672$, while the matched altered-connectivity control broadened to event distance $0.246737213404$ and robustness distance $0.03287037037$. Phase XVII then froze a thin causal slice on the dominant \texttt{bias\_onset} probe at $n_{\mathrm{side}} \in \{60,72,84\}$ and seeds $\{1303,1302\}$, obtaining mean edge reproducibility $0.626262626263$, acyclicity-band width $0.10101010101$, causal-depth drift $0.0$, and mismatch mean $0.03590611347$, while the control failed all matched gates. Phase XVIII finally reused only frozen Phase XV--XVII artifacts and showed that shell ordering, arrival-vs-depth scaling, and causal depth remain mutually compatible on the same authority slice, with branch max shell overlap $0.36577920989033486$, max adjacent slope drift $0.00476189285714268$, causal-depth drift $0.0$, and triangle-violation rate $0.0$, while the control fails both shell-ordering and refinement-scaling gates. The bounded conclusion is that the frozen branch supports an emergent-kinematics feasibility stack: ordering, causal closure, and a proto-geometric distance surrogate are mutually compatible. No claim is made about metric emergence, curvature, spacetime, continuum geometry, or physical correspondence.
\end{abstract}

\section{Scope and Frozen Contracts}

The synthesis covers only the already frozen Phase XVI--XVIII regime:
\begin{itemize}
\item Phase XVI: temporal-ordering feasibility;
\item Phase XVII: causal-closure feasibility;
\item Phase XVIII: proto-geometric distance-surrogate feasibility.
\end{itemize}

Nothing earlier is modified. The Phase V recovery readout, the Phase VI operator hierarchy $\Delta_h$, the Phase VII spectral-feasibility bundle, the Phase VIII short-time trace window, the Phase IX coefficient-stabilization bundle, the Phase X continuum-bridge bundle, the Phase XI survival-basin bundle, the Phase XII interaction bundle, the Phase XIII multi-mode sector bundle, and the Phase XIV collective-dynamics bundle all remain frozen inputs.

The concrete authority references are:
\begin{itemize}
\item \url{phase15-propagation/phase15_manifest.json};
\item \url{phase16-temporal-ordering/phase16_manifest.json};
\item \url{phase17-causal-closure/phase17_manifest.json};
\item \url{phase18-distance-surrogate/phase18_manifest.json}.
\end{itemize}

This paper does not introduce a new operator, a new disturbance family, a new seed set, or a new refinement hierarchy. It only synthesizes the existing frozen layer.

\section{Frozen Authority Slice}

The three-phase chain narrows to one common authority slice at the point where causal and distance claims become sharpest:
\begin{itemize}
\item disturbance family: \texttt{bias\_onset};
\item refinements: $n_{\mathrm{side}} \in \{60,72,84\}$;
\item ensemble size: $7$;
\item branch seeds: $\{1303,1302\}$;
\item control hierarchy: \texttt{periodic\_diagonal\_augmented\_control}.
\end{itemize}

Phase XVI uses a broader disturbance and ensemble regime to establish the ordering layer, but the primary nontrivial ordering probe is already \texttt{bias\_onset}. Phases XVII and XVIII then freeze that same slice into the causal and surrogate authority path.

\section{Phase XVI: Ordering Layer}

Phase XVI asks whether the collective sector supports a stable internal ordering relation built only from observable changes. The event vocabulary is frozen to threshold crossings of disturbance radius, dispersion, spectral response, low-$k$ response, width shift, and near/far front arrival. These events define ordered chains
\[
A \prec B \quad \Leftrightarrow \quad t_A < t_B,
\]
with deterministic lexical tie-breaking only for reproducibility.

On the primary \texttt{bias\_onset} probe, the branch metrics are:
\[
\text{primary event distance} = 0.142857142857,
\]
\[
\text{primary ordering score} = 0.965986394558,
\]
\[
\text{front-order distance} = 0.192592592593,
\]
\[
\text{ordering robustness distance} = 0.01421957672.
\]

The matched control broadens to:
\[
\text{primary event distance} = 0.246737213404,
\]
\[
\text{primary ordering score} = 0.941253044428,
\]
\[
\text{front-order distance} = 0.226543209877,
\]
\[
\text{ordering robustness distance} = 0.03287037037.
\]

Both hierarchies retain a monotonic cumulative evolution parameter, so monotonicity alone is not the separator. The separator is the full chain of event stability, front ordering, and robustness under weak perturbations.

\section{Phase XVII: Causal-Closure Layer}

Phase XVII is an intentionally thin authority probe. It consumes only frozen Phase XV--XVI artifacts, reconstructs influence edges from the event-chain ordering, and asks whether the resulting directed graph remains compact across the same refinement slice.

The branch aggregate metrics are:
\[
\text{mean edge reproducibility} = 0.626262626263,
\]
\[
\text{max adjacent edge-set distance} = 0.333333333333,
\]
\[
\text{acyclicity-band width} = 0.10101010101,
\]
\[
\text{mean causal depth} = 1.714285714286,
\qquad
\text{causal-depth drift} = 0.0,
\]
\[
\text{mismatch mean} = 0.03590611347.
\]

The control broadens materially:
\[
\text{mean edge reproducibility} = 0.585858585859,
\]
\[
\text{max adjacent edge-set distance} = 0.571428571429,
\]
\[
\text{acyclicity-band width} = 0.393939393939,
\]
\[
\text{causal-depth drift} = 0.714285714286,
\]
\[
\text{mismatch mean} = 0.037018301472.
\]

The bounded reading is narrow but important. Phase XVII does not prove physical causality. It shows that the branch supports a compact directed influence structure whose depth and ordering remain controlled while the matched control fails the same gates.

\section{Phase XVIII: Distance-Surrogate Layer}

Phase XVIII performs no new dynamics. It reuses only frozen Phase XV--XVII ledgers and defines
\[
\tilde d(i,j)=\operatorname{depth}(i,j)\,\bar v_{\mathrm{eff}}(h),
\qquad
\tilde \tau(i,j)=t_{\mathrm{arr}}(i,j),
\]
with the question restricted to monotone compatibility, not equality.

The branch clears all four surrogate gates:
\begin{itemize}
\item monotonic shell fraction: $1.0$;
\item max shell overlap: $0.36577920989033486$;
\item max adjacent slope drift: $0.00476189285714268$;
\item max causal-depth drift: $0.0$;
\item mean order compatibility: $0.96409388653$;
\item triangle-violation rate: $0.0$.
\end{itemize}

The control fails the matched shell-ordering and refinement-scaling gates:
\begin{itemize}
\item max shell overlap: $0.44471648285594$;
\item max adjacent slope drift: $0.5846153749999999$;
\item max causal-depth drift: $0.7142857142850001$.
\end{itemize}

The branch therefore preserves the same ordering/depth structure across the full authority slice in a way that the control does not.

\section{Integrated Chain}

Taken together, the three phases define one bounded refinement-stable chain.

\subsection*{Step 1: Ordering}
Observable threshold crossings can be arranged into stable branch event chains.

\subsection*{Step 2: Causal Closure}
Those same chains define a compact directed influence graph with stable mean depth and bounded acyclicity spread.

\subsection*{Step 3: Distance Surrogate}
The frozen causal depth and frozen propagation scaling remain compatible enough to support shell monotonicity and stable arrival-vs-depth scaling on the same authority slice.

The claim is not that a metric has emerged. The claim is only that the three levels are mutually consistent:
\[
\text{ordering}
\;\Rightarrow\;
\text{causal DAG slice}
\;\Rightarrow\;
\text{depth-based surrogate compatibility}.
\]

\section{Bounded Interpretation}

This synthesis supports an emergent-kinematics feasibility statement only in the minimal sense used here. The branch hierarchy admits:
\begin{itemize}
\item a stable ordering layer;
\item a compact causal-closure layer;
\item a compatible depth-based distance-surrogate layer.
\end{itemize}

It does not follow that a continuum metric exists. It does not follow that spacetime has emerged. It does not follow that curvature, causal cones, or geometric invariants are available. Those remain open questions for later, separately contracted phases.

\section{Conclusion}

The frozen Phase XVI--XVIII stack supports a bounded emergent-kinematics statement: ordering, causal closure, and proto-geometric distance-surrogate compatibility are jointly feasible on the frozen branch-local cochain-Laplacian hierarchy and jointly weaker on the matched deterministic control.

\end{document}
